\documentclass[11pt]{article}
\usepackage[margin=1in]{geometry}
\usepackage{amsmath,amssymb,amsthm,mathtools}
\usepackage{graphicx}
\usepackage{microtype}
\usepackage[colorlinks=true,linkcolor=blue,citecolor=blue,urlcolor=blue]{hyperref}

\newtheorem{theorem}{Theorem}
\newtheorem{lemma}{Lemma}
\newtheorem{remark}{Remark}
\newcommand{\Z}{\mathbb Z}
\newcommand{\N}{\mathbb N}
\newcommand{\rad}{\operatorname{rad}}
\newcommand{\ee}{\mathrm e}

\title{Nilpotents of Every Index for White's Lacunary Weight}
\author{Run \texttt{fa\_banach\_001}, lane 0
  (\texttt{agent\_lane\_00})}
\date{17 August 2026}

\begin{document}
\maketitle

\begin{abstract}
White constructs the weight
\(
 \omega(n)=\exp(|n|_{S_0})
\)
on \(\Z\), where \(S_0=\{2^{i^2}:i\geq0\}\), and asks whether the
radical of \(\ell^1(\Z,\omega)^{**}\) contains nilpotent elements of
arbitrarily high index.  We answer yes.  In fact, for every \(m\geq2\)
we construct a weak-star cluster point \(\Phi_m\) of normalized point
masses such that \(\Phi_m\) has nilpotency index exactly \(m\).  More
strongly, the principal left ideal
\(\ell^1(\Z,\omega)^{**}\Box\Phi_m\) has nilpotency index exactly \(m\).
The construction uses sparse balanced digits in the mixed-radix chain
\(2^{i^2}\mid2^{(i+1)^2}\): fewer than \(m\) copies create no carries,
whereas \(m\) copies create an independent word-length saving at every
common digit.
\end{abstract}

\noindent\fbox{\parbox{0.94\textwidth}{
\textbf{Status: candidate full result, likely valid.}
The literal open question following Theorem 1.3 of arXiv:1708.09635 is
resolved affirmatively for the paper's exact weight.  Novelty is plausible
but not certified; human review should focus on the balanced-digit lemma and
the passage from the uniform convolution estimate to the Arens product.}}

\section{Exact source question}

Let
\[
 q_i=2^{i^2}\quad(i\geq0),\qquad
 S_0=\{q_i:i\geq0\},\qquad
 \eta(n)=|n|_{S_0},\qquad
 \omega(n)=\ee^{\eta(n)}.
\]
Here \(\eta(n)\) is the least number of signed elements of \(S_0\)
whose sum is \(n\).  White proves that the radical of
\(A^{**}\), where \(A=\ell^1(\Z,\omega)\) and \(A^{**}\) carries the
first Arens product \(\Box\), contains a nonnilpotent element.  He then
asks \cite[p.~2]{white2018}:

\begin{center}
\includegraphics[width=0.94\textwidth]{figures/open_problem_crop.png}
\end{center}

\begin{quote}
``However, we do not know whether or not this example has nilpotent
elements of arbitrarily high index.''
\end{quote}

We prove the stronger statement that the radical contains a principal left
ideal of every prescribed finite nilpotency index.

\section{Idea of the proof}

\subsection*{Proof Intuition}

The generators \(q_i\) form a mixed-radix system with enormous even radices
\[
 R_i=\frac{q_{i+1}}{q_i}=2^{2i+1}.
\]
Fix the desired index \(m\).  At sparse positions \(r_j\), put a digit
\(d_j\) just larger than \(R_{r_j}/(2m)\), and let \(t_k\) be the integer
formed from the first \(k\) such digits.  If fewer than \(m\) prefixes are
added, every resulting digit remains strictly below half its radix, so the
word lengths add exactly.  With \(m\) prefixes, every position common to all
of them crosses the half-radix threshold.  Replacing \(m d_j\) copies of
\(q_{r_j}\) by one \(q_{r_j+1}\) and a small negative digit saves at least
one generator.  Sparsity keeps these carries disjoint.  Hence the normalized
product of \(m\) point masses has norm at most
\(\ee^{-\min(k_1,\ldots,k_m)}\), while every product of fewer than \(m\)
has norm one.

A weak-star cluster point \(\Phi_m\) therefore satisfies
\(\Phi_m^m=0\) but \(\Phi_m^{m-1}\neq0\).  The same estimate survives after
arbitrary bounded factors are inserted between the point masses: at the
finite level convolution on \(\Z\) is commutative, so those factors can be
grouped separately.  Consequently the entire left ideal
\(A^{**}\Box\Phi_m\) vanishes at the chosen power \(m\), placing
\(\Phi_m\) in the Jacobson radical.

\section{A balanced mixed-radix lemma}

We use the following elementary exact description of the word metric.

\begin{lemma}[strict balanced digits]\label{lem:balanced}
Suppose a finite expansion
\[
 x=\sum_{i\geq0}c_iq_i
\]
satisfies
\[
 |c_i|\leq \frac{R_i}{2}-1\qquad(i\geq0).
\]
Then
\[
 \eta(x)=\sum_{i\geq0}|c_i|.
\]
\end{lemma}

\begin{proof}
The displayed expansion gives the upper bound.  Consider any other finite
integer-coefficient representation \(x=\sum_i a_iq_i\).  Normalize it from
low positions to high positions.  Once positions below \(i\) agree with the
digits \(c_0,\ldots,c_{i-1}\), reduction modulo
\(q_{i+1}=R_iq_i\) gives
\(a_i=c_i+R_ih\) for some \(h\in\Z\).  Replace
\[
 (a_i,a_{i+1})\quad\hbox{by}\quad(c_i,a_{i+1}+h).
\]
This preserves the represented integer.  It does not increase the sum of
absolute coefficients, because
\[
 |c_i|+|h|\leq |c_i+R_ih|.
\]
Indeed, this is immediate when \(c_i h\geq0\).  When \(c_i h<0\), we have
\[
 |c_i+R_ih|-|c_i|-|h|
 =(R_i-1)|h|-2|c_i|
 \geq (R_i-1)-(R_i-2)>0.
\]
Iterating through the finitely many nonzero positions converts the arbitrary
representation into the displayed one without increasing its length.  Thus
the displayed length is minimal.
\end{proof}

\section{Full resolution}

\begin{theorem}\label{thm:main}
For the exact weight \(\omega(n)=\exp(|n|_{S_0})\) constructed in
\cite{white2018}, and for every integer \(m\geq2\), there exists
\[
 \Phi_m\in\rad\bigl(\ell^1(\Z,\omega)^{**},\Box\bigr)
\]
whose nilpotency index is exactly \(m\).  More strongly, the principal left
ideal
\[
 I_m=\ell^1(\Z,\omega)^{**}\Box\Phi_m
\]
has nilpotency index exactly \(m\).
\end{theorem}

\begin{proof}
Fix \(m\geq2\).  Choose \(N\geq2\) so large that
\(R_i\geq4m^2\) for every \(i\geq N\), and set
\[
 r_j=N+2(j-1),\qquad
 d_j=\left\lfloor\frac{R_{r_j}}{2m}\right\rfloor+1
 \quad(j\geq1).
\]
The choice of \(N\) gives
\begin{equation}\label{eq:thresholds}
 (m-1)d_j\leq \frac{R_{r_j}}2-1,
 \qquad
 m d_j\geq \frac{R_{r_j}}2+1.
\end{equation}
Indeed, the second inequality is immediate, while
\[
 (m-1)d_j
 \leq \frac{(m-1)R_{r_j}}{2m}+m-1
 =\frac{R_{r_j}}2-\frac{R_{r_j}}{2m}+m-1
 \leq\frac{R_{r_j}}2-1.
\]
Define
\[
 t_k=\sum_{j=1}^k d_jq_{r_j}\qquad(k\geq1).
\]
Lemma \ref{lem:balanced} and \eqref{eq:thresholds} show that
\begin{equation}\label{eq:length-tk}
 \eta(t_k)=\sum_{j=1}^kd_j.
\end{equation}

Take \(1\leq s\leq m\) and positive integers \(k_1,\ldots,k_s\).  In
the sum \(t_{k_1}+\cdots+t_{k_s}\), the coefficient at \(q_{r_j}\) is
\(h_jd_j\), where
\[
 h_j=|\{\ell:j\leq k_\ell\}|.
\]
If \(s<m\), then \(h_j\leq m-1\), so all digits are strictly balanced by
\eqref{eq:thresholds}.  Lemma \ref{lem:balanced} therefore gives
\begin{equation}\label{eq:no-carry}
 \eta(t_{k_1}+\cdots+t_{k_s})
 =\sum_{\ell=1}^s\eta(t_{k_\ell})
 \qquad(s<m).
\end{equation}

Now let \(s=m\), and write \(K=\min(k_1,\ldots,k_m)\).  For every
\(j\leq K\), replace the digit \(m d_j\) by
\[
 m d_jq_{r_j}
 =(m d_j-R_{r_j})q_{r_j}+q_{r_j+1}.
\]
The new digit at \(r_j\) has absolute value at most
\(R_{r_j}/2-1\), by the second inequality in
\eqref{eq:thresholds}; the digit at \(r_j+1\) is \(1\).  The positions
\(r_j\) are separated by two, so no carry position is another selected
position.  All remaining digits have multiplicity at most \(m-1\).
Thus the resulting expansion is strictly balanced, and Lemma
\ref{lem:balanced} gives the exact identity
\begin{equation}\label{eq:carry}
 \eta(t_{k_1}+\cdots+t_{k_m})
 =\sum_{\ell=1}^m\eta(t_{k_\ell})
  -\sum_{j=1}^{K}\sigma_j,
 \qquad
 \sigma_j=2m d_j-R_{r_j}-1\geq1.
\end{equation}

Put \(e_t=\delta_t/\omega(t)\).  Then \(\|e_t\|_\omega=1\), and
\begin{equation}\label{eq:norm}
 \|e_{u_1}*\cdots*e_{u_s}\|_\omega
 =\exp\!\left(
   \eta(u_1+\cdots+u_s)-\sum_{\ell=1}^s\eta(u_\ell)
 \right).
\end{equation}
Equations \eqref{eq:no-carry}--\eqref{eq:carry} imply
\begin{align}
 \|e_{t_{k_1}}*\cdots*e_{t_{k_s}}\|_\omega&=1
 &&(s<m),\label{eq:lower}\\
 \|e_{t_{k_1}}*\cdots*e_{t_{k_m}}\|_\omega&leq \ee^{-K}
 &&(s=m).\label{eq:upper}
\end{align}

Choose a weak-star convergent subnet of \((e_{t_k})\) in the unit ball of
\(A^{**}\), and call its limit \(\Phi_m\).  Since the weight \(\omega\)
itself is a norm-one functional in
\(A^*=\ell^\infty(\Z,1/\omega)\), equation \eqref{eq:lower} and the
iterated-limit definition of \(\Box\) give
\[
 \langle\Phi_m^{\Box s},\omega\rangle=1
 \qquad(1\leq s<m).
\]
In particular \(\Phi_m^{\Box(m-1)}\neq0\).  On the other hand,
\eqref{eq:upper} tends to zero uniformly whenever all \(k_\ell\) tend to
infinity.  Every iterated weak-star limit defining
\(\Phi_m^{\Box m}\) is therefore zero, so
\[
 \Phi_m^{\Box m}=0.
\]
Thus \(\Phi_m\) has nilpotency index exactly \(m\).

It remains to prove radical membership.  Let
\(\Psi_1,\ldots,\Psi_m\in A^{**}\), and choose bounded nets
\((a_{i,\alpha})\subset A\) converging weak-star to \(\Psi_i\).  Expand
\[
 (\Psi_1\Box\Phi_m)\Box\cdots\Box
 (\Psi_m\Box\Phi_m)
\]
using the iterated-limit definition of the first Arens product.  At every
finite stage the approximating product is, because convolution on \(\Z\)
is commutative,
\[
 (a_{1,\alpha_1}*\cdots*a_{m,\alpha_m})*
 (e_{t_{k_1}}*\cdots*e_{t_{k_m}}).
\]
Its norm is at most
\[
 \left(\prod_{i=1}^m\|a_{i,\alpha_i}\|_\omega\right)\ee^{-K}
\]
by \eqref{eq:upper}.  The nets \((a_{i,\alpha})\) are bounded and the
estimate is uniform in their indices, so every required iterated limit is
zero.  Hence \(I_m^{\Box m}=\{0\}\).  Conversely,
\(\Phi_m\in I_m\) because \(A^{**}\) is unital, and
\(\Phi_m^{\Box(m-1)}\neq0\), so \(I_m\) has index exactly \(m\).

A nilpotent left ideal consists of quasinilpotent elements and is contained
in the Jacobson radical.  Therefore \(\Phi_m\in\rad(A^{**},\Box)\), as
claimed.
\end{proof}

\section{Verification and boundary checks}

The proof uses strict inequalities, so there is no ambiguous half-radix
digit.  The spacing \(r_{j+1}-r_j=2\) is essential: it ensures that the
carry \(q_{r_j+1}\) does not alter another chosen digit.  The construction
also uses prefixes rather than disjoint blocks; consequently exactly the
first \(\min(k_1,\ldots,k_m)\) digits carry, producing the uniform decay
needed for every ordering of Arens limits.

The accompanying script independently minimizes finite word-length problems
by mixed-integer linear programming.  It checks all threshold formulas for
\(m=2,\ldots,8\), 32 exact word-length identities, and 69,632 instances of
the local normalization inequality.  These checks are corroborative only;
Lemma \ref{lem:balanced} supplies the proof for all \(m\).

\section{Novelty audit and limitations}

The local run indexes, exact-title and exact-question web searches, phrase
variants involving weighted \(\ell^1(\Z)\), the source author's current
publication list, and the author's 2026 sequel arXiv:2602.02764
\cite{white2026} were checked
through 17 August 2026.  The sequel concerns different radicals for the two
Arens products on a weighted free group and does not state this result.  The
OpenAlex record for the 2018 source reported no indexed citing works at the
time of the audit.  No prior answer to this exact question was located.
This is a bounded search, not a priority claim.

The theorem concerns the particular lacunary weight in the question and the
first Arens product, exactly as asked.  It does not settle the source's much
broader open questions about \(\rad(\ell^1(\Z)^{**})^2\), semisimplicity for
nonamenable groups, or Fourier algebras.

\begin{thebibliography}{9}
\bibitem{white2018}
J.~T. White,
\emph{The radical of the bidual of a Beurling algebra},
Quarterly Journal of Mathematics \textbf{69} (2018), 975--993;
arXiv:1708.09635.

\bibitem{white2026}
J.~T. White,
\emph{Radicals of Biduals of Beurling Algebras Can Be Different for the Two
Arens Products}, arXiv:2602.02764 (2026).
\end{thebibliography}

\end{document}
